\section{Produktutviklingsprosessen}
\label{sec:pu-prosess}
% Beskriv valgt utviklingsmetode med fagkilder, og hvordan prosessen er tilpasset prosjektet.
% Vis gjerne prosessen som en figur, med faser fra behov og krav via konsepter til prototype og testing.
\subsection{Iterativ prototyping}
\label{sec:iterativ-prototyping}
Etter at konseptet er valgt, utvikles bilen gjennom iterativ prototyping. Ifølge \textcite{camburn_2017} er en prototype en foreløpig modell som brukes før produksjon for å undersøke én eller flere sider ved et konsept eller design. Den trenger ikke være fysisk eller omfatte hele produktet. En CAD-modell, en simulering eller en prøvebit som bare tester én pasning regnes også som en prototype.

Fremgangsmåten følger \textcite{espedal_2024}. Hver del utvikles i versjoner, der den første er utgangspunktet. Hver versjon har én hovedtest som avgjør om endringen er en forbedring, og noen tilleggstester. Testene vurderes som vellykket, delvis vellykket eller mislykket og kobles til kravene i seksjon~\ref{sec:produktkrav}. Den beste versjonen så langt er gjeldende beste versjon og grunnlaget for neste iterasjon. Der det er praktisk, endres bare én parameter om gangen, slik at virkningen av endringen kan skilles fra andre forhold. Til slutt sammenlignes alle versjonene av en del i en samlet tabell.

Tabell~\ref{tab:iterasjonsserier} viser delene som utvikles på denne måten.


\begin{table}[H]
    \centering
    \caption{Iterasjonsserier for bilen}
    \label{tab:iterasjonsserier}
    \small
    \renewcommand{\arraystretch}{1.3}
    \begin{tabular}{|p{3.6cm}|p{6.0cm}|l|}
        \hline
        \textbf{Del} & \textbf{Hovedtest} & \textbf{Krav} \\
        \hline
        Drivverk (i CAD) & Geometrisk kontroll og simulert distanse & R1, R3 \\
        \hline
        Lagersete & Pasning mellom lager og sete & R8 \\
        \hline
        Feste av nav på aksel & Navet slurer ikke under opptrekk & R6, R7 \\
        \hline
        Snortrommel & Turtall etter opptrekk & R3, R4 \\
        \hline
        Drivhjul og dekk & Kjøredistanse & R3, R4, R5 \\
        \hline
        Ramme (generativ design) & Lukket ramme rundt drivverket & R1, R7, R8 \\
        \hline
    \end{tabular}
\end{table}

\subsubsection{Simulering som hjelpeverktøy}
\label{sec:simulering-verktoy}
I tillegg til de fysiske prøvebitene er det brukt en enkel numerisk simulering av én kjøring: bilen settes ned med svinghjulet i gang, hjulene slurer til de får tak, og bilen ruller til den stopper. Simuleringen er bare brukt til å sammenligne alternativer før de bygges, som utvekslingen i konseptvalget og versjonene av drivverket i CAD (seksjon~\ref{sec:dv-serien}). Flere parametere, blant annet rullemotstanden og virkningsgraden i tannhjulsinngrepet, er anslått, så de absolutte distansene er usikre. Rekkefølgen mellom alternativene er mindre følsom for dette. Utledning og validering av modellen hører til MAST2007.
